% Q13_problems.tex — the written problems of Q13 (Practice II-A). Built by
% `make -C paper paper` with Q13_answers.tex. B1 (a specification's soundness
% and completeness, a concept question) and B2 (the paper card) are the
% module's written pair: 15 minutes, after the chapter's §2 and before H3v.
\documentclass[11pt]{article}
% preamble.tex — shared preamble for the dafny-verus-learn-0 written track.
%
% Usage (every <ID>_problems.tex, <ID>_answers.tex and solutions/<ID>_solutions.tex):
%
%   \documentclass[11pt]{article}
%   % preamble.tex — shared preamble for the dafny-verus-learn-0 written track.
%
% Usage (every <ID>_problems.tex, <ID>_answers.tex and solutions/<ID>_solutions.tex):
%
%   \documentclass[11pt]{article}
%   % preamble.tex — shared preamble for the dafny-verus-learn-0 written track.
%
% Usage (every <ID>_problems.tex, <ID>_answers.tex and solutions/<ID>_solutions.tex):
%
%   \documentclass[11pt]{article}
%   \input{preamble}          % the Makefile always runs latexmk from paper/
%   \begin{document}
%   ...
%   \end{document}
%
% Engine: pdflatex. Dafny and Verus sources are mostly ASCII; the `literate`
% tables below map the common Unicode symbols to math glyphs inside listings.
% For other Unicode, extend a table, or build that file with lualatex
% (`latexmk -lualatex <file>`) after replacing the inputenc/fontenc lines with
% \usepackage{fontspec} and a monospace font with good coverage.
%
% Contents: program-logic macros (Hoare triples, wp, VCs, substitution,
% semantics brackets, typing judgments), inference rules (mathpartir),
% theorem and problem environments, the answer environment, the grader
% margin-note macro, and listings languages `dafny` and `verus` (Python uses
% listings' built-in `Python`), with inline forms \dfy{...}, \vrs{...}, \py{...}.
\usepackage[utf8]{inputenc}
\usepackage[T1]{fontenc}

\usepackage{amsmath}
\usepackage{amssymb}
\usepackage{amsthm}
\usepackage{stmaryrd}         % \llbracket \rrbracket for semantics brackets
\usepackage{mathpartir}       % inference rules: \inferrule, the mathpar environment
\usepackage{listings}
\usepackage{xcolor}
\usepackage[margin=1.1in]{geometry}

% ---------------------------------------------------------------------------
% Program logic.
%   \hoare{P}{c}{Q}       {P} c {Q}           (partial correctness)
%   \thoare{P}{c}{Q}      [P] c [Q]           (total correctness)
%   \provable{P}{c}{Q}    |- {P} c {Q}        (derivable in the proof system)
%   \WP{c}{Q}             wp(c, Q)            (\wp is LaTeX's Weierstrass p, so
%   \WLP{c}{Q}            wlp(c, Q)            the macro names are upper case)
%   \VC{c}{Q}             VC(c, Q)
%   \subst{Q}{x}{e}       Q[e/x]              (e replaces x in Q)
%   \sem{e}               [[e]]               (denotation)
%   \assign, \seqc        :=  and the sequencing semicolon
%   \impl                 ==> as a math arrow (LaTeX's own \iff gives <==>)
% ---------------------------------------------------------------------------
\newcommand{\hoare}[3]{\{\,#1\,\}\;#2\;\{\,#3\,\}}
\newcommand{\thoare}[3]{[\,#1\,]\;#2\;[\,#3\,]}
\newcommand{\provable}[3]{\vdash \hoare{#1}{#2}{#3}}
\newcommand{\WP}[2]{\mathrm{wp}(#1,\,#2)}
\newcommand{\WLP}[2]{\mathrm{wlp}(#1,\,#2)}
\newcommand{\VC}[2]{\mathrm{VC}(#1,\,#2)}
\newcommand{\subst}[3]{#1[#3/#2]}
\newcommand{\sem}[1]{\llbracket #1 \rrbracket}
\newcommand{\assign}{\mathrel{:=}}
\newcommand{\seqc}{\mathbin{;}\,}
\newcommand{\impl}{\Longrightarrow}

% ---------------------------------------------------------------------------
% Judgments.  \judgment{\Gamma}{e}{\tau}  ==>  Gamma |- e : tau
% ---------------------------------------------------------------------------
\newcommand{\judgment}[3]{#1 \vdash #2 : #3}
\newcommand{\judg}[3]{\judgment{#1}{#2}{#3}}          % short alias
\newcommand{\emptyjudgment}[2]{\cdot \vdash #1 : #2}  % empty context
\newcommand{\ctx}{\Gamma}
\newcommand{\turnstile}{\vdash}

% ---------------------------------------------------------------------------
% Theorem environments.
% ---------------------------------------------------------------------------
\theoremstyle{plain}
\newtheorem{theorem}{Theorem}[section]
\newtheorem{lemma}[theorem]{Lemma}
\newtheorem{corollary}[theorem]{Corollary}
\newtheorem{proposition}[theorem]{Proposition}
\theoremstyle{definition}
\newtheorem{definition}[theorem]{Definition}
% Written-track problems are numbered B1, B2, ... on both tracks (P<n> is a
% predict item). Label each one \label{prob:B<n>} and its answer \label{ans:B<n>}.
\newtheorem{problem}{Problem}
\renewcommand{\theproblem}{B\arabic{problem}}
\theoremstyle{remark}
\newtheorem*{remark}{Remark}
\newenvironment{answer}
  {\par\medskip\noindent\textbf{Answer.}\begin{quotation}\noindent}
  {\end{quotation}\medskip}

% ---------------------------------------------------------------------------
% Grader margin notes. Rendered form of the `% GRADER[...][...][...]:` comment
% convention (see template_answers.tex). Use inline where a visible note helps:
%   \gradernote{B2/induction-soundness/2: the invariant is not preserved by the last assignment}
% ---------------------------------------------------------------------------
\definecolor{gradercolor}{RGB}{160, 30, 30}
\newcommand{\gradernote}[1]{\textcolor{gradercolor}{\small\sffamily [GRADER: #1]}}

% ---------------------------------------------------------------------------
% Code listings: Dafny, Verus (Rust inside verus!{...}), Python.
% ---------------------------------------------------------------------------
\definecolor{codekeyword}{RGB}{0, 60, 130}
\definecolor{codecomment}{RGB}{90, 110, 90}
\definecolor{codestring}{RGB}{130, 60, 10}
\definecolor{codehole}{RGB}{180, 0, 0}

\lstdefinelanguage{dafny}{
  sensitive=true,
  morekeywords={
    abstract, allocated, as, assert, assume, bool, break, by, calc, case, char,
    class, codatatype, const, constructor, continue, datatype, decreases, else,
    ensures, exists, expect, export, extends, false, forall, fresh, function,
    ghost, if, import, in, include, int, invariant, is, iterator, label, lemma,
    map, match, method, modifies, modify, module, multiset, nat, new, newtype,
    null, old, opaque, opened, predicate, print, provides, reads, real, refines,
    requires, return, returns, reveal, reveals, seq, set, static, string, then,
    this, trait, true, twostate, type, unchanged, var, while, witness},
  emph={assume},
  emphstyle=\color{codehole}\bfseries,
  comment=[l]{//},
  morecomment=[s]{/*}{*/},
  morestring=[b]",
}

\lstdefinelanguage{verus}{
  sensitive=true,
  morekeywords={
    as, break, const, continue, crate, dyn, else, enum, false, fn, for, if,
    impl, in, let, loop, match, mod, move, mut, pub, ref, return, self, Self,
    static, struct, super, trait, true, type, unsafe, use, where, while,
    spec, proof, exec, open, closed, tracked, ghost, requires, ensures,
    recommends, invariant, decreases, assert, forall, exists, choose, by,
    broadcast, uninterp, axiom, returns, int, nat, Seq, Set, Map, Ghost, Tracked},
  emph={assume, admit},
  emphstyle=\color{codehole}\bfseries,
  comment=[l]{//},
  morecomment=[s]{/*}{*/},
  morestring=[b]",
}

\lstset{
  basicstyle=\ttfamily\small,
  keywordstyle=\color{codekeyword}\bfseries,
  commentstyle=\color{codecomment}\itshape,
  stringstyle=\color{codestring},
  showstringspaces=false,
  columns=fullflexible,
  keepspaces=true,
  breaklines=true,
  frame=leftline,
  framerule=0.6pt,
  xleftmargin=1em,
  aboveskip=0.8em,
  belowskip=0.8em,
  literate=
    {→}{{$\to$}}1        {←}{{$\leftarrow$}}1  {⇒}{{$\Rightarrow$}}1
    {⇔}{{$\Leftrightarrow$}}1                  {↦}{{$\mapsto$}}1
    {∀}{{$\forall$}}1    {∃}{{$\exists$}}1     {λ}{{$\lambda$}}1
    {∧}{{$\wedge$}}1     {∨}{{$\vee$}}1        {¬}{{$\neg$}}1
    {≠}{{$\neq$}}1       {≤}{{$\leq$}}1        {≥}{{$\geq$}}1
    {∈}{{$\in$}}1        {∉}{{$\notin$}}1      {⊆}{{$\subseteq$}}1
    {×}{{$\times$}}1     {·}{{$\cdot$}}1       {∅}{{$\varnothing$}}1
    {ℕ}{{$\mathbb{N}$}}1 {ℤ}{{$\mathbb{Z}$}}1
    {⊢}{{$\vdash$}}1     {⟦}{{$\llbracket$}}1  {⟧}{{$\rrbracket$}}1,
}

% Code environments and inline code.
\lstnewenvironment{dafnycode}{\lstset{language=dafny}}{}
\lstnewenvironment{veruscode}{\lstset{language=verus}}{}
\lstnewenvironment{pycode}{\lstset{language=Python}}{}
\newcommand{\dfy}[1]{\lstinline[language=dafny]{#1}}
\newcommand{\vrs}[1]{\lstinline[language=verus]{#1}}
\newcommand{\py}[1]{\lstinline[language=Python]{#1}}
          % the Makefile always runs latexmk from paper/
%   \begin{document}
%   ...
%   \end{document}
%
% Engine: pdflatex. Dafny and Verus sources are mostly ASCII; the `literate`
% tables below map the common Unicode symbols to math glyphs inside listings.
% For other Unicode, extend a table, or build that file with lualatex
% (`latexmk -lualatex <file>`) after replacing the inputenc/fontenc lines with
% \usepackage{fontspec} and a monospace font with good coverage.
%
% Contents: program-logic macros (Hoare triples, wp, VCs, substitution,
% semantics brackets, typing judgments), inference rules (mathpartir),
% theorem and problem environments, the answer environment, the grader
% margin-note macro, and listings languages `dafny` and `verus` (Python uses
% listings' built-in `Python`), with inline forms \dfy{...}, \vrs{...}, \py{...}.
\usepackage[utf8]{inputenc}
\usepackage[T1]{fontenc}

\usepackage{amsmath}
\usepackage{amssymb}
\usepackage{amsthm}
\usepackage{stmaryrd}         % \llbracket \rrbracket for semantics brackets
\usepackage{mathpartir}       % inference rules: \inferrule, the mathpar environment
\usepackage{listings}
\usepackage{xcolor}
\usepackage[margin=1.1in]{geometry}

% ---------------------------------------------------------------------------
% Program logic.
%   \hoare{P}{c}{Q}       {P} c {Q}           (partial correctness)
%   \thoare{P}{c}{Q}      [P] c [Q]           (total correctness)
%   \provable{P}{c}{Q}    |- {P} c {Q}        (derivable in the proof system)
%   \WP{c}{Q}             wp(c, Q)            (\wp is LaTeX's Weierstrass p, so
%   \WLP{c}{Q}            wlp(c, Q)            the macro names are upper case)
%   \VC{c}{Q}             VC(c, Q)
%   \subst{Q}{x}{e}       Q[e/x]              (e replaces x in Q)
%   \sem{e}               [[e]]               (denotation)
%   \assign, \seqc        :=  and the sequencing semicolon
%   \impl                 ==> as a math arrow (LaTeX's own \iff gives <==>)
% ---------------------------------------------------------------------------
\newcommand{\hoare}[3]{\{\,#1\,\}\;#2\;\{\,#3\,\}}
\newcommand{\thoare}[3]{[\,#1\,]\;#2\;[\,#3\,]}
\newcommand{\provable}[3]{\vdash \hoare{#1}{#2}{#3}}
\newcommand{\WP}[2]{\mathrm{wp}(#1,\,#2)}
\newcommand{\WLP}[2]{\mathrm{wlp}(#1,\,#2)}
\newcommand{\VC}[2]{\mathrm{VC}(#1,\,#2)}
\newcommand{\subst}[3]{#1[#3/#2]}
\newcommand{\sem}[1]{\llbracket #1 \rrbracket}
\newcommand{\assign}{\mathrel{:=}}
\newcommand{\seqc}{\mathbin{;}\,}
\newcommand{\impl}{\Longrightarrow}

% ---------------------------------------------------------------------------
% Judgments.  \judgment{\Gamma}{e}{\tau}  ==>  Gamma |- e : tau
% ---------------------------------------------------------------------------
\newcommand{\judgment}[3]{#1 \vdash #2 : #3}
\newcommand{\judg}[3]{\judgment{#1}{#2}{#3}}          % short alias
\newcommand{\emptyjudgment}[2]{\cdot \vdash #1 : #2}  % empty context
\newcommand{\ctx}{\Gamma}
\newcommand{\turnstile}{\vdash}

% ---------------------------------------------------------------------------
% Theorem environments.
% ---------------------------------------------------------------------------
\theoremstyle{plain}
\newtheorem{theorem}{Theorem}[section]
\newtheorem{lemma}[theorem]{Lemma}
\newtheorem{corollary}[theorem]{Corollary}
\newtheorem{proposition}[theorem]{Proposition}
\theoremstyle{definition}
\newtheorem{definition}[theorem]{Definition}
% Written-track problems are numbered B1, B2, ... on both tracks (P<n> is a
% predict item). Label each one \label{prob:B<n>} and its answer \label{ans:B<n>}.
\newtheorem{problem}{Problem}
\renewcommand{\theproblem}{B\arabic{problem}}
\theoremstyle{remark}
\newtheorem*{remark}{Remark}
\newenvironment{answer}
  {\par\medskip\noindent\textbf{Answer.}\begin{quotation}\noindent}
  {\end{quotation}\medskip}

% ---------------------------------------------------------------------------
% Grader margin notes. Rendered form of the `% GRADER[...][...][...]:` comment
% convention (see template_answers.tex). Use inline where a visible note helps:
%   \gradernote{B2/induction-soundness/2: the invariant is not preserved by the last assignment}
% ---------------------------------------------------------------------------
\definecolor{gradercolor}{RGB}{160, 30, 30}
\newcommand{\gradernote}[1]{\textcolor{gradercolor}{\small\sffamily [GRADER: #1]}}

% ---------------------------------------------------------------------------
% Code listings: Dafny, Verus (Rust inside verus!{...}), Python.
% ---------------------------------------------------------------------------
\definecolor{codekeyword}{RGB}{0, 60, 130}
\definecolor{codecomment}{RGB}{90, 110, 90}
\definecolor{codestring}{RGB}{130, 60, 10}
\definecolor{codehole}{RGB}{180, 0, 0}

\lstdefinelanguage{dafny}{
  sensitive=true,
  morekeywords={
    abstract, allocated, as, assert, assume, bool, break, by, calc, case, char,
    class, codatatype, const, constructor, continue, datatype, decreases, else,
    ensures, exists, expect, export, extends, false, forall, fresh, function,
    ghost, if, import, in, include, int, invariant, is, iterator, label, lemma,
    map, match, method, modifies, modify, module, multiset, nat, new, newtype,
    null, old, opaque, opened, predicate, print, provides, reads, real, refines,
    requires, return, returns, reveal, reveals, seq, set, static, string, then,
    this, trait, true, twostate, type, unchanged, var, while, witness},
  emph={assume},
  emphstyle=\color{codehole}\bfseries,
  comment=[l]{//},
  morecomment=[s]{/*}{*/},
  morestring=[b]",
}

\lstdefinelanguage{verus}{
  sensitive=true,
  morekeywords={
    as, break, const, continue, crate, dyn, else, enum, false, fn, for, if,
    impl, in, let, loop, match, mod, move, mut, pub, ref, return, self, Self,
    static, struct, super, trait, true, type, unsafe, use, where, while,
    spec, proof, exec, open, closed, tracked, ghost, requires, ensures,
    recommends, invariant, decreases, assert, forall, exists, choose, by,
    broadcast, uninterp, axiom, returns, int, nat, Seq, Set, Map, Ghost, Tracked},
  emph={assume, admit},
  emphstyle=\color{codehole}\bfseries,
  comment=[l]{//},
  morecomment=[s]{/*}{*/},
  morestring=[b]",
}

\lstset{
  basicstyle=\ttfamily\small,
  keywordstyle=\color{codekeyword}\bfseries,
  commentstyle=\color{codecomment}\itshape,
  stringstyle=\color{codestring},
  showstringspaces=false,
  columns=fullflexible,
  keepspaces=true,
  breaklines=true,
  frame=leftline,
  framerule=0.6pt,
  xleftmargin=1em,
  aboveskip=0.8em,
  belowskip=0.8em,
  literate=
    {→}{{$\to$}}1        {←}{{$\leftarrow$}}1  {⇒}{{$\Rightarrow$}}1
    {⇔}{{$\Leftrightarrow$}}1                  {↦}{{$\mapsto$}}1
    {∀}{{$\forall$}}1    {∃}{{$\exists$}}1     {λ}{{$\lambda$}}1
    {∧}{{$\wedge$}}1     {∨}{{$\vee$}}1        {¬}{{$\neg$}}1
    {≠}{{$\neq$}}1       {≤}{{$\leq$}}1        {≥}{{$\geq$}}1
    {∈}{{$\in$}}1        {∉}{{$\notin$}}1      {⊆}{{$\subseteq$}}1
    {×}{{$\times$}}1     {·}{{$\cdot$}}1       {∅}{{$\varnothing$}}1
    {ℕ}{{$\mathbb{N}$}}1 {ℤ}{{$\mathbb{Z}$}}1
    {⊢}{{$\vdash$}}1     {⟦}{{$\llbracket$}}1  {⟧}{{$\rrbracket$}}1,
}

% Code environments and inline code.
\lstnewenvironment{dafnycode}{\lstset{language=dafny}}{}
\lstnewenvironment{veruscode}{\lstset{language=verus}}{}
\lstnewenvironment{pycode}{\lstset{language=Python}}{}
\newcommand{\dfy}[1]{\lstinline[language=dafny]{#1}}
\newcommand{\vrs}[1]{\lstinline[language=verus]{#1}}
\newcommand{\py}[1]{\lstinline[language=Python]{#1}}
          % the Makefile always runs latexmk from paper/
%   \begin{document}
%   ...
%   \end{document}
%
% Engine: pdflatex. Dafny and Verus sources are mostly ASCII; the `literate`
% tables below map the common Unicode symbols to math glyphs inside listings.
% For other Unicode, extend a table, or build that file with lualatex
% (`latexmk -lualatex <file>`) after replacing the inputenc/fontenc lines with
% \usepackage{fontspec} and a monospace font with good coverage.
%
% Contents: program-logic macros (Hoare triples, wp, VCs, substitution,
% semantics brackets, typing judgments), inference rules (mathpartir),
% theorem and problem environments, the answer environment, the grader
% margin-note macro, and listings languages `dafny` and `verus` (Python uses
% listings' built-in `Python`), with inline forms \dfy{...}, \vrs{...}, \py{...}.
\usepackage[utf8]{inputenc}
\usepackage[T1]{fontenc}

\usepackage{amsmath}
\usepackage{amssymb}
\usepackage{amsthm}
\usepackage{stmaryrd}         % \llbracket \rrbracket for semantics brackets
\usepackage{mathpartir}       % inference rules: \inferrule, the mathpar environment
\usepackage{listings}
\usepackage{xcolor}
\usepackage[margin=1.1in]{geometry}

% ---------------------------------------------------------------------------
% Program logic.
%   \hoare{P}{c}{Q}       {P} c {Q}           (partial correctness)
%   \thoare{P}{c}{Q}      [P] c [Q]           (total correctness)
%   \provable{P}{c}{Q}    |- {P} c {Q}        (derivable in the proof system)
%   \WP{c}{Q}             wp(c, Q)            (\wp is LaTeX's Weierstrass p, so
%   \WLP{c}{Q}            wlp(c, Q)            the macro names are upper case)
%   \VC{c}{Q}             VC(c, Q)
%   \subst{Q}{x}{e}       Q[e/x]              (e replaces x in Q)
%   \sem{e}               [[e]]               (denotation)
%   \assign, \seqc        :=  and the sequencing semicolon
%   \impl                 ==> as a math arrow (LaTeX's own \iff gives <==>)
% ---------------------------------------------------------------------------
\newcommand{\hoare}[3]{\{\,#1\,\}\;#2\;\{\,#3\,\}}
\newcommand{\thoare}[3]{[\,#1\,]\;#2\;[\,#3\,]}
\newcommand{\provable}[3]{\vdash \hoare{#1}{#2}{#3}}
\newcommand{\WP}[2]{\mathrm{wp}(#1,\,#2)}
\newcommand{\WLP}[2]{\mathrm{wlp}(#1,\,#2)}
\newcommand{\VC}[2]{\mathrm{VC}(#1,\,#2)}
\newcommand{\subst}[3]{#1[#3/#2]}
\newcommand{\sem}[1]{\llbracket #1 \rrbracket}
\newcommand{\assign}{\mathrel{:=}}
\newcommand{\seqc}{\mathbin{;}\,}
\newcommand{\impl}{\Longrightarrow}

% ---------------------------------------------------------------------------
% Judgments.  \judgment{\Gamma}{e}{\tau}  ==>  Gamma |- e : tau
% ---------------------------------------------------------------------------
\newcommand{\judgment}[3]{#1 \vdash #2 : #3}
\newcommand{\judg}[3]{\judgment{#1}{#2}{#3}}          % short alias
\newcommand{\emptyjudgment}[2]{\cdot \vdash #1 : #2}  % empty context
\newcommand{\ctx}{\Gamma}
\newcommand{\turnstile}{\vdash}

% ---------------------------------------------------------------------------
% Theorem environments.
% ---------------------------------------------------------------------------
\theoremstyle{plain}
\newtheorem{theorem}{Theorem}[section]
\newtheorem{lemma}[theorem]{Lemma}
\newtheorem{corollary}[theorem]{Corollary}
\newtheorem{proposition}[theorem]{Proposition}
\theoremstyle{definition}
\newtheorem{definition}[theorem]{Definition}
% Written-track problems are numbered B1, B2, ... on both tracks (P<n> is a
% predict item). Label each one \label{prob:B<n>} and its answer \label{ans:B<n>}.
\newtheorem{problem}{Problem}
\renewcommand{\theproblem}{B\arabic{problem}}
\theoremstyle{remark}
\newtheorem*{remark}{Remark}
\newenvironment{answer}
  {\par\medskip\noindent\textbf{Answer.}\begin{quotation}\noindent}
  {\end{quotation}\medskip}

% ---------------------------------------------------------------------------
% Grader margin notes. Rendered form of the `% GRADER[...][...][...]:` comment
% convention (see template_answers.tex). Use inline where a visible note helps:
%   \gradernote{B2/induction-soundness/2: the invariant is not preserved by the last assignment}
% ---------------------------------------------------------------------------
\definecolor{gradercolor}{RGB}{160, 30, 30}
\newcommand{\gradernote}[1]{\textcolor{gradercolor}{\small\sffamily [GRADER: #1]}}

% ---------------------------------------------------------------------------
% Code listings: Dafny, Verus (Rust inside verus!{...}), Python.
% ---------------------------------------------------------------------------
\definecolor{codekeyword}{RGB}{0, 60, 130}
\definecolor{codecomment}{RGB}{90, 110, 90}
\definecolor{codestring}{RGB}{130, 60, 10}
\definecolor{codehole}{RGB}{180, 0, 0}

\lstdefinelanguage{dafny}{
  sensitive=true,
  morekeywords={
    abstract, allocated, as, assert, assume, bool, break, by, calc, case, char,
    class, codatatype, const, constructor, continue, datatype, decreases, else,
    ensures, exists, expect, export, extends, false, forall, fresh, function,
    ghost, if, import, in, include, int, invariant, is, iterator, label, lemma,
    map, match, method, modifies, modify, module, multiset, nat, new, newtype,
    null, old, opaque, opened, predicate, print, provides, reads, real, refines,
    requires, return, returns, reveal, reveals, seq, set, static, string, then,
    this, trait, true, twostate, type, unchanged, var, while, witness},
  emph={assume},
  emphstyle=\color{codehole}\bfseries,
  comment=[l]{//},
  morecomment=[s]{/*}{*/},
  morestring=[b]",
}

\lstdefinelanguage{verus}{
  sensitive=true,
  morekeywords={
    as, break, const, continue, crate, dyn, else, enum, false, fn, for, if,
    impl, in, let, loop, match, mod, move, mut, pub, ref, return, self, Self,
    static, struct, super, trait, true, type, unsafe, use, where, while,
    spec, proof, exec, open, closed, tracked, ghost, requires, ensures,
    recommends, invariant, decreases, assert, forall, exists, choose, by,
    broadcast, uninterp, axiom, returns, int, nat, Seq, Set, Map, Ghost, Tracked},
  emph={assume, admit},
  emphstyle=\color{codehole}\bfseries,
  comment=[l]{//},
  morecomment=[s]{/*}{*/},
  morestring=[b]",
}

\lstset{
  basicstyle=\ttfamily\small,
  keywordstyle=\color{codekeyword}\bfseries,
  commentstyle=\color{codecomment}\itshape,
  stringstyle=\color{codestring},
  showstringspaces=false,
  columns=fullflexible,
  keepspaces=true,
  breaklines=true,
  frame=leftline,
  framerule=0.6pt,
  xleftmargin=1em,
  aboveskip=0.8em,
  belowskip=0.8em,
  literate=
    {→}{{$\to$}}1        {←}{{$\leftarrow$}}1  {⇒}{{$\Rightarrow$}}1
    {⇔}{{$\Leftrightarrow$}}1                  {↦}{{$\mapsto$}}1
    {∀}{{$\forall$}}1    {∃}{{$\exists$}}1     {λ}{{$\lambda$}}1
    {∧}{{$\wedge$}}1     {∨}{{$\vee$}}1        {¬}{{$\neg$}}1
    {≠}{{$\neq$}}1       {≤}{{$\leq$}}1        {≥}{{$\geq$}}1
    {∈}{{$\in$}}1        {∉}{{$\notin$}}1      {⊆}{{$\subseteq$}}1
    {×}{{$\times$}}1     {·}{{$\cdot$}}1       {∅}{{$\varnothing$}}1
    {ℕ}{{$\mathbb{N}$}}1 {ℤ}{{$\mathbb{Z}$}}1
    {⊢}{{$\vdash$}}1     {⟦}{{$\llbracket$}}1  {⟧}{{$\rrbracket$}}1,
}

% Code environments and inline code.
\lstnewenvironment{dafnycode}{\lstset{language=dafny}}{}
\lstnewenvironment{veruscode}{\lstset{language=verus}}{}
\lstnewenvironment{pycode}{\lstset{language=Python}}{}
\newcommand{\dfy}[1]{\lstinline[language=dafny]{#1}}
\newcommand{\vrs}[1]{\lstinline[language=verus]{#1}}
\newcommand{\py}[1]{\lstinline[language=Python]{#1}}


\title{Q13 --- Practice II-A \\ \large Written problems}
\author{dafny-verus-learn-0}
\date{}

\begin{document}
\maketitle

\section*{Instructions}
Answer in \texttt{Q13\_answers.tex}, one section per problem; keep its
\verb+\label{ans:B<n>}+ anchors. \texttt{make -C paper paper} must stay green:
a non-compiling answer set is the written track's red build.

Order: after the chapter, before H3v, the prose item (15 minutes in all).
Answer B1 without running anything: the grader runs the checks. Then run
\texttt{/grade-paper Q13} in a fresh session. Compare with
\texttt{paper/solutions/} only as the module's last step.

% ---------------------------------------------------------------------------
\begin{problem}[concept]\label{prob:B1}
\textbf{Sound, complete, both or neither} (5 minutes). The function
\dfy{Ramp(x: int): int} returns $0$ for a negative \dfy{x} and \dfy{x}
otherwise; its reference postcondition is
\[ Q(x, y) \;\equiv\; y = (\textbf{if } x < 0 \textbf{ then } 0 \textbf{ else } x). \]
Card 05.5 calls a candidate postcondition $\hat Q$ \emph{sound} when
$\hat Q(x, y) \Rightarrow Q(x, y)$ for all $x, y$ (every output it allows is
the reference's), and \emph{complete} when $Q(x, y) \Rightarrow \hat Q(x, y)$
for all $x, y$ (it allows every reference output); the precondition is
\textsf{true}. Five candidates:
\begin{enumerate}
\item[$Q_1$] $y \ge 0$
\item[$Q_2$] $y \ge 0 \wedge y \ge x$
\item[$Q_3$] $y = x \vee y = 0$
\item[$Q_4$] $(x < 0 \Rightarrow y = 0) \wedge (x \ge 0 \Rightarrow y = x)$
\item[$Q_5$] $y = x$
\end{enumerate}
For each: \textbf{(a)} sound? complete? \textbf{(b)} for every property it
lacks, a concrete pair $(x, y)$ that shows it. \textbf{(c)} In one line: which
of the five would a client of \dfy{Ramp} prefer, and which would its
implementer find easiest to prove, and why are those different?
\end{problem}

% ---------------------------------------------------------------------------
\begin{problem}[paper-card]\label{prob:B2}
\textbf{Paper card 05.5, VERINA: Benchmarking Verifiable Code Generation}
(10 minutes). Fill the card on its eight fields: task shape, oracle, what is
trusted, cheat surface, metric, budget, variance, artifact. Each field is
graded against the pinned range of your own \texttt{cheatsheets.md}: lines
2005--2031, the entry from its heading to its \emph{Why it matters}
paragraph, in the copy the command below names. The pin is the sha256 of
exactly those lines, computed with the first line below; it must print the
second.
{\small
\begin{verbatim}
sed -n '2005,2031p' ~/code/research/agentic-formal-veri/cheatsheets.md | sha256sum
6c9b00d83e072910f2c5c6a2a298646045a89ffae04a90bdc0821436413f67f0  -
\end{verbatim}
}
Write each field in your own words, a sentence or two. Where the range says
nothing about a field, write that, and what you would need to know. A hash
that no longer matches means the card file changed: the grader refuses the
card until the pin is renewed.
\end{problem}

\end{document}
